\section{Sample Final Examination}

\subsubsection*{Problem 1 (20 points)}

\term{Question.} Let $S_1(t)$ and $S_2(t)$ be two asset price processes. Under the risk-neutral measure $\mathbb Q$, their stochastic differential equations are

$\dfrac{dS_1(t)}{S_1(t)}=r\,dt+\sigma_1\,dW_1(t),\qquad
\dfrac{dS_2(t)}{S_2(t)}=r\,dt+\sigma_2\,dW_2(t)$,

where $r$ is the continuously compounded risk-free annual interest rate, $\sigma_1$ and $\sigma_2$ are the volatilities of $S_1(t)$ and $S_2(t)$ respectively, and $W_1(t)$ and $W_2(t)$ are two independent standard Brownian motions under $\mathbb Q$.

Consider a $T$-year European option with payoff

$\operatorname{Payoff}=\max(S_1(T)-S_2(T),0)$.

Let $V$ be the price of the option at time 0. Using the likelihood ratio method,

$\dfrac{\partial^2V}{\partial S_1(0)\partial S_2(0)}$
is found to be in the form $\mathbb E^{\mathbb Q}[f(Z_1,Z_2)]$, where $\mathbb E^{\mathbb Q}[\cdot]$ is the expectation under the risk-neutral measure $\mathbb Q$, and $Z_1,Z_2$ are two independent standard normal random variables. Find $f(Z_1,Z_2)$.

\term{Solution.} From the given SDEs,

$S_i(t)=S_i(0)\exp\left[\left(r-\frac12\sigma_i^2\right)t+\sigma_iW_i(t)\right],\qquad i=1,2$.

Let $g(x_1,x_2)$ be the joint pdf of $S_1(t)$ and $S_2(t)$. Since they are independent,

$g(x_1,x_2)=g^{(1)}(x_1)g^{(2)}(x_2)$,
where $g^{(i)}(x_i)$, $i=1,2$, is the probability density function of $S_i(t)$,
$g^{(i)}(x_i)=\dfrac1{x_i\sigma_i\sqrt{2\pi T}}\exp\left[-\dfrac{\xi^2(x_i)}2\right]$,

with
$\xi(x_i)=\dfrac{\log(x_i/S_i(0))-(r-\frac12\sigma_i^2)T}{\sigma_i\sqrt T}$.

By the likelihood ratio method,

$\dfrac{\partial^2V}{\partial S_1(0)\partial S_2(0)}$
$=\mathbb E^{\mathbb Q}\!\left[
\dfrac{\frac{\partial^2g(S_1(T),S_2(T))}{\partial S_1(0)\partial S_2(0)}}{g(S_1(T),S_2(T))}
e^{-rT}\max(S_1(T)-S_2(T),0)\right]$.

Since $g(S_1(T),S_2(T))=g^{(1)}(S_1(T))g^{(2)}(S_2(T))$,

$\dfrac{\frac{\partial^2g(S_1(T),S_2(T))}{\partial S_1(0)\partial S_2(0)}}{g(S_1(T),S_2(T))}$
$=\dfrac{\frac{\partial g^{(1)}(S_1(T))}{\partial S_1(0)}}{g^{(1)}(S_1(T))}
\dfrac{\frac{\partial g^{(2)}(S_2(T))}{\partial S_2(0)}}{g^{(2)}(S_2(T))}$

$=\dfrac{\partial\log g^{(1)}(S_1(T))}{\partial S_1(0)}
\dfrac{\partial\log g^{(2)}(S_2(T))}{\partial S_2(0)}$.

For $i=1,2$,
$\dfrac{\partial\log g^{(i)}(S_i(T))}{\partial S_i(0)}
=-\xi(S_i(T))\dfrac{\partial\xi(S_i(T))}{\partial S_i(0)}
=\dfrac{\xi(S_i(T))}{S_i(0)\sigma_i\sqrt T}$.

Let $Z_1,Z_2$ be independent standard normal random variables. Then

$S_i(t)=S_i(0)\exp\left[\left(r-\frac12\sigma_i^2\right)t+\sigma_i\sqrt t\,Z_i\right]$,

so
$\dfrac{\partial\log g^{(i)}(S_i(T))}{\partial S_i(0)}
=\dfrac{Z_i}{S_i(0)\sigma_i\sqrt T}$.

Hence
$\dfrac{\frac{\partial^2g(S_1(T),S_2(T))}{\partial S_1(0)\partial S_2(0)}}{g(S_1(T),S_2(T))}
=\dfrac{Z_1Z_2}{S_1(0)S_2(0)\sigma_1\sigma_2T}$.

Therefore
$\dfrac{\partial^2V}{\partial S_1(0)\partial S_2(0)}
=\mathbb E^{\mathbb Q}[f(Z_1,Z_2)]$,

where
$f(Z_1,Z_2)=\dfrac{e^{-rT}Z_1Z_2}{S_1(0)S_2(0)\sigma_1\sigma_2T}
\max(S_1(T)-S_2(T),0)$,

with
$S_1(T)=S_1(0)\exp\left[\left(r-\frac12\sigma_1^2\right)T+\sigma_1\sqrt T Z_1\right]$,

$S_2(T)=S_2(0)\exp\left[\left(r-\frac12\sigma_2^2\right)T+\sigma_2\sqrt T Z_2\right]$.

\subsubsection*{Problem 2 (25 points)}

\term{Question.} Let $U$ be uniformly distributed on $[0,1]$. Estimate $\mathbb E[e^{2U}]$ by stratified sampling, partitioning $[0,1]$ into $k=3$ strata of equal length.

$\mathrm{(a)}$ With total sample size $3n$, write down the stratified estimator with proportional allocation.

$\mathrm{(b)}$ Evaluate the magnitude of variance reduction from the estimator in (a).

$\mathrm{(c)}$ Find the optimal allocation and write down the corresponding estimator.

\term{Solution.} $\mathrm{(a)}$ Let $Y=e^{2U}$ and partition $[0,1]$ into

$A_1=[0,1/3],\qquad A_2=(1/3,2/3],\qquad A_3=(2/3,1]$.

With proportional allocation, $p_i=1/3$ and $n_i=p_i(3n)=n$, $i=1,2,3$. The stratified estimator is

$\widehat Y=\dfrac1{3n}\sum_{i=1}^3\sum_{j=1}^nY_{ij}$,

where
$Y_{1j}=e^{2V_{1j}},\qquad Y_{2j}=e^{2V_{2j}},\qquad Y_{3j}=e^{2V_{3j}}$,

and $V_{ij}$, $j=1,\ldots,n$, are independent and uniformly distributed on $A_i$, $i=1,2,3$.

$\mathrm{(b)}$ To evaluate the magnitude of the variance reduction, we need

$\mu_1=\mathbb E[Y_{1j}]=\int_0^{1/3}3e^{2x}\,dx=\frac32(e^{2/3}-1)$,

$\mu_2=\mathbb E[Y_{2j}]=\int_{1/3}^{2/3}3e^{2x}\,dx=\frac32(e^{4/3}-e^{2/3})$,

$\mu_3=\mathbb E[Y_{3j}]=\int_{2/3}^{1}3e^{2x}\,dx=\frac32(e^2-e^{4/3})$.

Thus the magnitude of variance reduction is

$\dfrac1{3n}\left[\sum_{i=1}^3p_i\mu_i^2-\left(\sum_{i=1}^3p_i\mu_i\right)^2\right]$

$=\dfrac1{3n}\left(\dfrac13\right)\left(\dfrac94\right)
\left[(e^{2/3}-1)^2+(e^{4/3}-e^{2/3})^2+(e^2-e^{4/3})^2\right]$

$\quad-\dfrac1{3n}\left(\dfrac19\right)\left(\dfrac94\right)
\left[(e^{2/3}-1)+(e^{4/3}-e^{2/3})+(e^2-e^{4/3})\right]^2$
$=\dfrac{0.9065}{n}$.

$\mathrm{(c)}$ The optimal allocation is
$q_i^*=\dfrac{p_i\sigma_i}{\sum_{l=1}^3p_l\sigma_l}$,
where $\sigma_i^2=\operatorname{Var}(Y_{ij})$. 

For the first stratum,
$\operatorname{Var}(Y_{1j})=\operatorname{Var}(e^{2V_{1j}})$
$=\mathbb E[e^{4V_{1j}}]-\{\mathbb E[e^{2V_{1j}}]\}^2$

$=\int_0^{1/3}3e^{4x}\,dx-\left(\int_0^{1/3}3e^{2x}\,dx\right)^2$
$=\frac34(e^{4/3}-1)-\frac94(e^{2/3}-1)^2=0.0743$.

For the second stratum,

$\operatorname{Var}(Y_{2j})=\operatorname{Var}(e^{2V_{2j}})$
$=\int_{1/3}^{2/3}3e^{4x}\,dx-\left(\int_{1/3}^{2/3}3e^{2x}\,dx\right)^2$

$=\frac34(e^{8/3}-e^{4/3})-\frac94(e^{4/3}-e^{2/3})^2=0.2819$.

For the third stratum,
$\operatorname{Var}(Y_{3j})=\operatorname{Var}(e^{2V_{3j}})$
$=\int_{2/3}^{1}3e^{4x}\,dx-\left(\int_{2/3}^{1}3e^{2x}\,dx\right)^2$

$=\frac34(e^4-e^{8/3})-\frac94(e^2-e^{4/3})^2=1.0693$.

Therefore
$q_1^*=\dfrac{\frac13\sqrt{0.0743}}{\frac13(\sqrt{0.0743}+\sqrt{0.2819}+\sqrt{1.0693})}=0.1483$,

$q_2^*=\dfrac{\frac13\sqrt{0.2819}}{\frac13(\sqrt{0.0743}+\sqrt{0.2819}+\sqrt{1.0693})}=0.2889$,

$q_3^*=\dfrac{\frac13\sqrt{1.0693}}{\frac13(\sqrt{0.0743}+\sqrt{0.2819}+\sqrt{1.0693})}=0.5627$.

The stratified estimator corresponding to the optimal allocation is

$\widehat Y^*=\dfrac1{9n}\left(
\dfrac1{q_1^*}\sum_{j=1}^{3nq_1^*}Y_{1j}
+\dfrac1{q_2^*}\sum_{j=1}^{3nq_2^*}Y_{2j}
+\dfrac1{q_3^*}\sum_{j=1}^{3nq_3^*}Y_{3j}\right)$.

\subsubsection*{Problem 3 (15 points)}

\term{Question.} Suppose that we wish to generate samples from a probability distribution with density

$f(x)=\begin{cases}\dfrac{e^{-x}}{1-e^{-a}},&0\leq x\leq a,\\0,&\text{otherwise},\end{cases}$

by the acceptance-rejection method. The alternative density is

$g(x)=\begin{cases}\lambda e^{-\lambda x},&x\geq0,\\0,&\text{otherwise},\end{cases}$

for some $\lambda>0$. Determine the best $\lambda$ that minimizes the expected number of trial samples needed to generate a sample from $f(x)$.

\term{Solution.} Let

$h(x)=\dfrac{f(x)}{g(x)}=
\begin{cases}
\dfrac{e^{x(\lambda-1)}}{\lambda(1-e^{-a})},&0\leq x\leq a,\\
0,&\text{otherwise}.
\end{cases}$

For $0<\lambda<1$, $h(x)$ is monotonically decreasing and maximal at $x=0$, with maximal value
$h(\lambda)=\dfrac1{\lambda(1-e^{-a})}$.
For $\lambda=1$,
$h(\lambda)=h(1)=\dfrac1{1-e^{-a}}$.

For $\lambda>1$, $h(x)$ is monotonically increasing and maximal at $x=a$, with maximal value
$h(\lambda)=\dfrac{e^{a(\lambda-1)}}{\lambda(1-e^{-a})}$.
Under the acceptance-rejection method, the expected number of trial samples needed to generate a sample from $f$ is $h(\lambda)$. The minimal value for $h(\lambda)$ is attained at $\lambda=1$ and
$h(1)=\dfrac1{1-e^{-a}}$.

\subsubsection*{Problem 4 (20 points)}

\term{Question.} Let $\alpha>1$ and $\beta>0$. Suppose the probability density function of $X$ and $Y$ is

$g(x,y)\propto\begin{cases}(1-x)^{\beta-1}y^{\alpha-2},&0<y<x<1,\\0,&\text{otherwise}.
\end{cases}$

Construct a Gibbs sampling method to generate $X$ and $Y$ from $g(x,y)$. You must include the pseudocode of your Gibbs sampling method.

\term{Solution.} Construct a Markov chain $\{\mathbf x^{(M)}=(x^{(M)},y^{(M)})\}$ whose limiting distribution converges to $g(x,y)$. Let $g(x)$ and $g(y)$ be the marginal pdfs of $X$ and $Y$ respectively.

For $y<x<1$ and $0<y<1$, the conditional pdf of $X$ given $Y=y$ is

$g(x\mid y)=\dfrac{g(x,y)}{g(y)}$
$=\dfrac{(1-x)^{\beta-1}y^{\alpha-2}}{\int_y^1(1-x)^{\beta-1}y^{\alpha-2}\,dx}$
$=\dfrac{(1-x)^{\beta-1}}{(1-y)^\beta/\beta}
=\dfrac{\beta(1-x)^{\beta-1}}{(1-y)^\beta}$.

The conditional distribution function of $X$ given $Y=y$ is

$G(x\mid y)=\int_y^xg(s\mid y)\,ds$
$=\int_y^x\dfrac{\beta(1-s)^{\beta-1}}{(1-y)^\beta}\,ds$
$=\dfrac{\beta}{(1-y)^\beta}\left[-\dfrac{(1-s)^\beta}{\beta}\right]_y^x$

$=1-\dfrac{(1-x)^\beta}{(1-y)^\beta}$.
By the inverse transform method, given $Y=y$, the variate $X$ can be generated from a uniform random number $u$ as follows:
$u=1-\dfrac{(1-x)^\beta}{(1-y)^\beta}$,

$(1-x)^\beta=(1-y)^\beta(1-u)$,
$x=1-(1-y)(1-u)^{1/\beta}$.

Similarly, for $0<y<x$ and $0<x<1$, the conditional pdf of $Y$ given $X=x$ is

$g(y\mid x)=\dfrac{g(x,y)}{g(x)}$
$=\dfrac{(1-x)^{\beta-1}y^{\alpha-2}}{\int_0^x(1-x)^{\beta-1}y^{\alpha-2}\,dy}$
$=\dfrac{(\alpha-1)y^{\alpha-2}}{x^{\alpha-1}}$.

The conditional distribution function of $Y$ given $X=x$ is

$G(y\mid x)=\int_0^yg(s\mid x)\,ds$

$=\int_0^y\dfrac{(\alpha-1)s^{\alpha-2}}{x^{\alpha-1}}\,ds$
$=\dfrac{y^{\alpha-1}}{x^{\alpha-1}}$.

By the inverse transform method, given $X=x$, the variate $Y$ can be generated from a uniform random number $u$ as follows:
$u=\dfrac{y^{\alpha-1}}{x^{\alpha-1}},\qquad y=xu^{1/(\alpha-1)}$.

\term{Pseudocode.}

Step 1: Set $i=1$ and $y^{(0)}=a$, where $a$ is a specified constant in $(0,1)$.

Step 2: Generate $U\sim U(0,1)$ and set
$x^{(i)}=1-(1-y^{(i-1)})(1-U)^{1/\beta}$.

Step 3: Generate $V\sim U(0,1)$ and set
$y^{(i)}=x^{(i)}V^{1/(\alpha-1)}$.

Step 4: Set $i=i+1$.

Step 5: Go to Step 2 until $i$ equals a prespecified integer $M$.

\subsubsection*{Problem 5 (20 points)}

\term{Question.} For a 0.75-year American call option on a non-dividend paying stock, the strike is $\$68$, exercise is allowed only at 0.25, 0.5, and 0.75 year, and the continuously compounded risk-free annual interest rate is $4.8\%$. Let $S(t)$ be the stock price. Under the risk-neutral measure, 11 sample paths are

\noindent\scalebox{0.827}{$\begin{array}{c|rrr}
\text{Path}&t=0.25&t=0.5&t=0.75\\\hline
1&62.98&65.57&57.36\\2&85.06&68.26&48.36\\3&96.43&86.03&95.55\\
4&72.84&74.31&79.49\\5&73.53&69.11&68.79\\6&72.24&83.65&62.49\\
7&74.44&84.15&90.62\\8&71.35&68.15&60.27\\9&78.69&93.53&81.25\\
10&60.25&64.22&75.78\\11&85.09&70.59&66.31
\end{array}$}

Using the least squares approach, assume that the expected payoff from continuation at time $t_i=0.25i$, $i=1,2$, is modelled by a linear function of $S(t_i)$. Calculate the value of the American call option at time 0.

\term{Solution.} Terminal payoffs $Y_3$ are

\noindent\scalebox{0.827}{$\begin{array}{c|r|r}
\text{Path}&S(0.75)&Y_3=\max(S(0.75)-68,0)\\\hline
1&57.36&0\\2&48.36&0\\3&95.55&27.55\\4&79.49&11.49\\5&68.79&0.79\\
6&62.49&0\\7&90.62&22.62\\8&60.27&0\\9&81.25&13.25\\10&75.78&7.78\\11&66.31&0
\end{array}$}

At $t=0.5$, model the expected payoff from continuation as a linear function $f_2(S(0.5))$ of $S(0.5)$. Its coefficients are estimated from the in-the-money data by least squares:

\noindent\scalebox{0.827}{$\begin{array}{c|r|r|c}
\text{Path}&Y_3e^{-0.048(0.25)}=Y_3e^{-0.012}&S(0.5)&\text{Exercise in-the-money?}\\\hline
1&0&65.57&N\\2&0&68.26&Y\\3&27.22&86.03&Y\\4&11.35&74.31&Y\\
5&0.78&69.11&Y\\6&0&83.65&Y\\7&22.35&84.15&Y\\8&0&68.15&Y\\
9&13.09&93.53&Y\\10&7.69&64.22&N\\11&0&70.59&Y
\end{array}$}

Applying the regression function in the scientific calculator to the in-the-money data,

$\mathbb E[Y_3e^{-0.012}\mid S(0.5)]
=0.7653S(0.5)-51.0253=:f_2(S(0.5))$.

With this conditional expectation function, the payoffs $Y_2$ at $t=0.5$ are

$Y_2=\begin{cases}
S(0.5)-68,&S(0.5)-68\geq f_2(S(0.5)),\\
e^{-0.012}Y_3,&\text{otherwise}.
\end{cases}$

The following table gives $Y_2$ for each sample path:

\noindent\scalebox{0.827}{$\begin{array}{c|r|r|r|r}
\text{Path}&\text{Exercise }S(0.5)-68&\text{Continuation }f_2(S(0.5))&Y_3e^{-0.012}&Y_2\\\hline
1&0&0&0&0\\2&0.26&1.21&0&0\\3&18.03&14.81&27.22&18.03\\
4&6.31&5.84&11.35&6.31\\5&1.11&1.86&0.78&0.78\\6&15.65&12.99&0&15.65\\
7&16.15&13.37&22.35&16.15\\8&0.15&1.13&0&0\\9&25.53&20.55&13.09&25.53\\
10&0&0&7.69&7.69\\11&2.59&3&0&0
\end{array}$}

Since Paths 1 and 10 are out-of-money, $Y_2=Y_3e^{-0.012}$ on those paths.

Next, repeat the procedure for $t=0.25$ using

\noindent\scalebox{0.827}{$\begin{array}{c|r|r|c}
\text{Path}&Y_2e^{-0.012}&S(0.25)&\text{Exercise in-the-money?}\\\hline
1&0&62.98&N\\2&0&85.06&Y\\3&17.81&96.43&Y\\4&6.23&72.84&Y\\
5&0.77&73.53&Y\\6&15.46&72.24&Y\\7&15.96&74.44&Y\\8&0&71.35&Y\\
9&25.23&78.69&Y\\10&7.6&60.25&N\\11&0&85.09&Y
\end{array}$}

We have

$\mathbb E[Y_2e^{-0.012}\mid S(0.25)]
=0.1661S(0.25)-4.0499=:f_1(S(0.25))$.

With this conditional expectation function, the payoffs $Y_1$ at $t=0.25$ are

$Y_1=\begin{cases}
S(0.25)-68,&S(0.25)-68\geq f_1(S(0.25)),\\
e^{-0.012}Y_2,&\text{otherwise}.
\end{cases}$

The following table gives $Y_1$ for each sample path:

\noindent\scalebox{0.827}{$\begin{array}{c|r|r|r|r}
\text{Path}&\text{Exercise }S(0.25)-68&\text{Continuation }f_1(S(0.25))&Y_2e^{-0.012}&Y_1\\\hline
1&0&0&0&0\\2&17.06&10.08&0&17.06\\3&28.43&11.97&17.81&28.43\\
4&4.84&8.05&6.23&6.23\\5&5.53&8.16&0.77&0.77\\6&4.24&7.95&15.46&15.46\\
7&6.44&8.31&15.96&15.96\\8&3.35&7.8&0&0\\9&10.69&9.02&25.23&10.69\\
10&0&0&7.6&7.6\\11&17.09&10.08&0&17.09
\end{array}$}

Since Paths 1 and 10 are out-of-money, $Y_1=Y_2e^{-0.012}$ on those paths.

Finally, the price of the American option at $t=0$ is estimated by the average of $e^{-0.012}Y_1$, that is, $10.72$.
